\section{Método de integração e testes}\label{cap:testes}
A avaliação foi organizada em inspeção de projeto, testes automatizados, execução nas placas e integração de ponta a ponta. O desenvolvimento incluiu testes de caixa preta, exercitando entradas, saídas e estados observáveis, e de caixa branca, verificando regras internas, persistência e tratamento de erros. O procedimento distingue testes executados durante o desenvolvimento do protocolo controlado ainda a concluir.

\subsection{Preparação, configuração e unidade de ensaio}
Antes de cada bloco, devem ser registrados revisão mecânica, disco, unidade física de teste, alimentação, versões de firmware/runtime/API/Care, ambiente e política de travas. A referência temporal, a calibração, a identidade dos módulos e a ausência de operação incerta são condições de preparação. Instrumento, operador e responsável pela observação acompanham o registro.

Uma tentativa de dispensação corresponde a um pedido identificado e seu desfecho, inclusive quando recusado ou interrompido. ``Unidade de teste'' designa o objeto físico usado no ensaio, com tipo e dimensões identificados; não significa necessariamente medicamento destinado a ingestão. O número de pedidos, de ciclos internos, de unidades solicitadas e de unidades coletadas é conservado separadamente.

\subsection{Dispensação com contagem independente}
A campanha controlada consiste em repetir um roteiro sob condições registradas, mantendo constantes as características escolhidas para o bloco. A quantidade no copo é contada por observação física, sem adotar como verdade o estado do motor, a leitura do sensor ou a baixa no aplicativo. Essa é a referência independente. Sua comparação com o pedido permite identificar acerto, ausência, unidade extra, retenção, travamento e resultado não observado.

O procedimento compreende conferir o estoque inicial, identificar o pedido, executar sem reenvio indevido, contar o coletor, registrar tempo e condição do mecanismo e conferir os estoques físico e digital. Fotografias identificadas apoiam a auditoria. Ausência de observação não é registrada como zero. Mudança de peça, ajuste ou versão inicia novo bloco; não elimina falhas do anterior. Fichas e proposta de blocos estão no Apêndice~\ref{ap:ensaios}.

\subsection{Mecânica, elétrica e sensoriamento}
A inspeção dimensional compara as cotas nominais V75 às medidas da peça impressa, sem inferir tolerância apenas da malha. A avaliação elétrica considera polaridade, continuidade, tensão e corrente em repouso e atuação, quedas nos contatos e aquecimento. Remoção energizada exige ensaio próprio; reconexão lógica não comprova segurança elétrica.

Para o sensor, a referência é a inspeção da cavidade vazia e ocupada, com o reservatório montado. Localização, calibração e restauração de posição são avaliadas em condições identificadas, inclusive perda de referência. Resultados de ADC parado não são extrapolados para o mecanismo em movimento.

\subsection{Software, comunicação e atualização}
As suítes automatizadas verificam estados, contratos, identidade, quadros válidos e inválidos, CRC, duplicatas e persistência. No aplicativo, consideram componentes, navegação e API simulada; no servidor, regras, contratos e gravação. Os registros conservam comando de execução, versão, casos aprovados, falhas, erros e ignorados. Rodadas sobrepostas não são somadas.

Em placa, são considerados parada, perda de resposta, reinício durante movimento, resultado retido e reconciliação. Atualizações são avaliadas por transferência, verificação de integridade, ativação, reinício e recuperação. Para a Pi, verifica-se também assinatura; nas ESP32, SHA-256 não é apresentado como assinatura. Os testes com dependências simuladas são distinguidos da atuação física.

\subsection{Integração Care--API--central e recuperação}
O roteiro acompanha uma sessão autenticada desde a solicitação até comando, atuação, confirmação humana, estoque e histórico, utilizando a mesma identidade de operação. Perda de rede, resposta tardia, expiração de sessão, duplicata e reinício são tratados junto à recuperação. Na operação sem internet, é preciso verificar relógio, agenda e prontidão, e não apenas a existência de SQLite.

A configuração entregue utiliza um paciente por cadastro. A validação atual não exige criar outro usuário ou outro paciente. Ampliação para múltiplos contextos e testes de acesso cruzado pertencem às possibilidades de continuidade. Isso não dispensa verificar credenciais inválidas, sessão e permissões das funções existentes.

\subsection{Telemetria e análise de medidas}
A coleta passiva registra tentativas, consultas válidas, falhas HTTP, estado das placas, identificadores de inicialização, leituras ADC, temperatura e filas. Uma consulta pode reutilizar o mesmo estado de hardware; por isso, amostras e leituras distintas são contadas separadamente. Mediana, percentis, extremos e lacunas são interpretados com a condição de carga e os desligamentos intencionais. Latência HTTP local, reconhecimento de LED, resposta de voz e duração da dose não são reunidos como uma única métrica.

\subsection{Avaliação da visão e do agente}
Na visão, a referência humana identifica sequências simuladas completas, mão vazia, objeto sem aproximação e sequência interrompida, com luz, distância e oclusão registradas. Ajuste e avaliação devem usar conjuntos separados. A máquina de estados final é avaliada independentemente dos escores de versões anteriores.

Para precisão $P$, sensibilidade $R$ e acurácia, com VP, FP, VN e FN definidos pela referência, utilizam-se:
\begin{equation}\label{eq:vis-metricas}
P=\frac{VP}{VP+FP},\quad R=\frac{VP}{VP+FN},\quad
\mathrm{Acc}=\frac{VP+VN}{VP+VN+FP+FN},\quad
\mathrm{FPS}=\frac{N_q}{\Delta t}.
\end{equation}
Denominador nulo resulta em métrica não definida. Inconclusivos e falhas técnicas são reportados separadamente, sem interpretar ausência de detecção como prova de ausência de tomada.

No agente, verificam-se contexto do paciente único, seleção de ferramenta, argumentos, confirmação, recusa, efeito persistido, ambiguidades, falhas e limites de finalidade. Ferramentas MCP, resposta do modelo e experiência de voz são unidades de avaliação diferentes. As versões do agente são identificadas para que melhorias históricas não sejam atribuídas indistintamente à versão final.

\subsection{Rastreabilidade e análise}
A matriz relaciona requisito, caso, configuração, evidência e desfecho. Tentativas interrompidas, falhas e itens não avaliados permanecem identificados. Os resultados numéricos e sua interpretação são apresentados exclusivamente na Seção~\ref{cap:resultados}, sem antecipar resultados empíricos nesta seção metodológica.

\subsection*{\textcolor{DCPendente}{Complementações necessárias}}
\pendencia{TST01}{Inserir protocolo aprovado de bancada: instrumentos, revisão, condições, comprimidos de teste, repetições, referência independente, limites de parada e critérios de aceitação. Não usar amostra exploratória como campanha final.}
\pendencia{TST02}{Fornecer relatórios completos por suíte: comando de execução, versão, ambiente, casos, aprovados, falhas/erros, data e arquivos de evidência. Não somar contagens de revisões sobrepostas nem atribuí-las ao Care, à visão ou ao agente.}
\pendencia{TST03}{Executar e documentar a integração de ponta a ponta e os cenários sem rede/sem NTP, somente após comissionamento. Incluir RTC, reinício, reconexão, filas e correspondência entre eventos locais e centrais.}
\pendencia{TST04}{Completar os protocolos e relatórios de visão e agente com versões, dados, métricas e falhas. Distinguir servidor MCP, seleção de ferramentas, qualidade de resposta, voz e integração com o Care.}
\pendencia{TST05}{Consolidar a matriz final requisito--teste--evidência--resultado, incluindo itens não executados e falhas abertas. Vincular cada resultado aa Seção~\ref{cap:resultados}; não antecipar discussão ou conclusão.}
